\documentclass[10pt,a4paper]{amsart}
\usepackage[margin=19mm]{geometry}
\usepackage{amsmath,amssymb,mathtools}
\usepackage[scr=rsfs,cal=euler]{mathalfa}
\usepackage{xfrac}
\usepackage[unicode,hidelinks,bookmarks=false]{hyperref}
\newcommand{\ie}{i.e.\ }
\newcommand{\cf}{cf.\ }
\newcommand{\resp}{respectively\ }
\newcommand{\Subsection}{\subsection}
\providecommand{\nobreakdash}{\nobreak\hbox{-}}
\setlength{\emergencystretch}{2em}
\multlinegap0pt
\usepackage{amssymb}
\usepackage{mathtools}
\usepackage{enumerate}
\newtheorem{corollary}{Corollary}
\newtheorem{proposition}{Proposition}
\newcommand{\Z}{\mathbb{Z}}
\title{The geodesic restriction problem for arithmetic spherical harmonics}
\author{Maximiliano Sanchez Garza}
\date{Source: arXiv:2509.24874v2 (CC BY 4.0)}
\begin{document}
\maketitle

\noindent\emph{Source and licence.} The geodesic restriction problem for arithmetic spherical harmonics. Version arXiv:2509.24874v2 (CC BY 4.0), distributed by arXiv under the Creative Commons Attribution 4.0 International licence, http://creativecommons.org/licenses/by/4.0/, as recorded in the arXiv licence metadata for this exact version and shown on the abstract page https://arxiv.org/abs/2509.24874v2. Full licence notice: \texttt{licenses/arxiv-2509.24874-CC-BY-4.0.txt}.\par\smallskip

\noindent\textit{Editorial scope.} Counted unit: the article's proposition BoundInnerSumsSa (source0.tex lines 1023--1030). The article's complete original proof of it is delivered with it (source0.tex lines 1032--1061, one \texttt{proof} environment of 30 lines). No mathematics is written, repaired or re-derived: the statement and the proof are the article's own lines, in the article's own notation, and no retained line is substituted or re-flowed.  Retained uncounted as the source-local context the proof needs: lines 1017--1019.  Nothing was omitted from any retained span, so the package carries no omission record.  No retained line contains a citation, so the wrapper carries no bibliography and no bibliography entry.  No retained line contains a figure environment, an image-inclusion command or drawing code, so no image is delivered and the package declares no asset dependency.  Editorial formatting only, in addition: an AMS \texttt{amsart} preamble that restates the article's own declarations and macros used by the retained lines, namely the environments \texttt{corollary}, \texttt{proposition} on separate counters, together with the standard packages those lines need.  The retained environments print in this document as Corollary 1, Proposition 1 in order of appearance, whereas the article prints the same items with the numbering of its own full document.  The complete native source of the article as distributed in the arXiv e-print for this exact version is bundled unchanged as \texttt{sources/185982/source0.tex}.
\begin{corollary}\label{BoundHATransform}
For any $j\geq 0$ and $\xi\neq 0$, we have $$ \widehat{H_A}(\alpha,m,\xi) \ll_{D,\varepsilon,j} \frac{T \ell^{\varepsilon}}{|T\xi|^{j}} (\ell TA)^{-1/2}.$$ Moreover, for $N(\alpha) \notin [A/2,2A]$, we have $\widehat{H_A}(\alpha,m,\xi)=0$.
\end{corollary}

\begin{proposition}\label{BoundInnerSumsSa}
For any $m\geq 1$ and $0 < \delta \leq 1$, we have 
\begin{align*}
    &\sum_{\alpha \in \mathcal{F}_{D}} \lambda_{f_{\psi}}(N(\alpha)) \sum_{\xi\in \Z} \widehat{H_A}\left( \alpha,m,\xi-\frac{\mathrm{Arg}(\alpha/\overline{\alpha})}{2\pi} \right) \\
    &\ll_{D,\varepsilon,\delta} \ell^{\varepsilon} T^{1/2} (\ell A)^{-1/2} \sum_{\alpha \in S_{T^{-1+\delta}}} |\lambda_{f_{\psi}}(N(\alpha))| + \ell^{\varepsilon} T^{-3/2} (\ell A)^{-1/2} \sum_{\substack{\alpha \in \mathcal{F}_{D} \\ A/2 \leq N(\alpha) \leq 2A}} |\lambda_{f_{\psi}}(N(\alpha))|,
\end{align*}
where $$ S_{R} := \left\{ \alpha \in \mathcal{F}_{D} : -\pi R \leq \mathrm{Arg}(\alpha) \leq \pi R, \: A/2 \leq N(\alpha) \leq 2A \right\}$$ for any $R\in (0,1)$.
\end{proposition}

\begin{proof}
From the definition of $\mathcal{F}_{D}$, for any $\alpha \in \mathcal{F}_{D}$ we have that $-\pi/2 \leq \mathrm{Arg}(\alpha) < \pi/2$, so that $-\pi \leq \mathrm{Arg}(\alpha/\overline{\alpha}) < \pi$. This implies that $-\frac{1}{2} \leq \frac{\mathrm{Arg}(\alpha / \overline{\alpha})}{2\pi} < \frac{1}{2}$ for any $\alpha \in \mathcal{F}_{D}$. Thus, from Corollary \ref{BoundHATransform} with $j=2$ and the triangle inequality, we have that
\begin{align}
    &\sum_{\alpha \in \mathcal{F}_{D}} \lambda_{f_{\psi}}(N(\alpha)) \sum_{\xi\geq 1} \widehat{H_A}\left( \alpha,m,\xi-\frac{\mathrm{Arg}(\alpha/\overline{\alpha})}{2 \pi} \right) \label{BoundInnerSumsSa-1} \\
    &\ll_{D,\varepsilon} \ell^{\varepsilon} T^{-3/2} (\ell A)^{-1/2} \sum_{\substack{\alpha \in \mathcal{F}_{D} \\ A/2 \leq N(\alpha) \leq 2A}} |\lambda_{f_{\psi}}(N(\alpha))| \sum_{\xi\geq 1} \frac{1}{\left| \xi - \frac{\mathrm{Arg}(\alpha/\overline{\alpha})}{2 \pi} \right|^2} \nonumber \\
    &\ll_{D,\varepsilon} \ell^{\varepsilon} T^{-3/2} (\ell A)^{-1/2} \sum_{\substack{\alpha \in \mathcal{F}_{D} \\ A/2 \leq N(\alpha) \leq 2A}} |\lambda_{f_{\psi}}(N(\alpha))| \sum_{\xi\geq 1} \frac{1}{\xi^2} \nonumber \\
    &\ll_{D,\varepsilon} \ell^{\varepsilon} T^{-3/2} (\ell A)^{-1/2} \sum_{\substack{\alpha \in \mathcal{F}_{D} \\ A/2 \leq N(\alpha) \leq 2A}} |\lambda_{f_{\psi}}(N(\alpha))|. \nonumber
\end{align}
Using the same reasoning, we have
\begin{equation}
    \sum_{\alpha \in \mathcal{F}_{D}} \lambda_{f_{\psi}}(N(\alpha)) \sum_{\xi\leq -1} \widehat{H_A}\left( \alpha,m,\xi-\frac{\mathrm{Arg}(\alpha/\overline{\alpha})}{2 \pi} \right) \ll_{D,\varepsilon} \ell^{\varepsilon} T^{-3/2} (\ell A)^{-1/2} \sum_{\substack{\alpha \in \mathcal{F}_{D} \\ A/2 \leq N(\alpha) \leq 2A}} |\lambda_{f_{\psi}}(N(\alpha))|. \label{BoundInnerSumsSa-2}
\end{equation}
Now, let $\xi = 0$ and $\alpha \in \mathcal{F}_{D}$. By the definition of $S_{R}$, we have that $\left|-\frac{\mathrm{Arg}(\alpha/\overline{\alpha})}{2 \pi} \right| \leq T^{-1+\delta}$ if and only if $\alpha \in S_{T^{-1+\delta}}$. If this is the case, from Corollary \ref{BoundHATransform} with $j=0$, we get 
\begin{equation}
    \widehat{H_A}\left(\alpha,m,-\frac{\mathrm{Arg}(\alpha/\overline{\alpha})}{2 \pi} \right) \ll_{D,\varepsilon} \ell^{\varepsilon} T^{1/2} (\ell A)^{-1/2}. \label{BoundInnerSumsSa-3}
\end{equation}
If not, then we have $\left|-\frac{\mathrm{Arg}(\alpha/\overline{\alpha})}{2 \pi} \right| > T^{-1+\delta}$. In particular, $\frac{\mathrm{Arg}(\alpha/\overline{\alpha})}{2 \pi} \neq 0$. Hence, again by Corollary \ref{BoundHATransform} with $j=\left\lceil 2/\delta \right\rceil$ we have 
\begin{align}
    \widehat{H_A}\left(\alpha,m,-\frac{\mathrm{Arg}(\alpha/\overline{\alpha})}{2\pi} \right) &\ll_{D,\varepsilon,\delta} \frac{\ell^{\varepsilon} T}{\left| T\left( -\frac{\mathrm{Arg}(\alpha/\overline{\alpha})}{\pi} \right) \right|^{j}}  (\ell TA)^{-1/2} \label{BoundInnerSumsSa-4} \\
    &\ll_{D,\varepsilon,\delta} \ell^{\varepsilon} T^{1-\delta j} (\ell TA)^{-1/2} \ll_{D,\varepsilon,\delta} \ell^{\varepsilon} T^{-3/2} (\ell A)^{-1/2}. \nonumber
\end{align}
Combining \eqref{BoundInnerSumsSa-3} and \eqref{BoundInnerSumsSa-4}, we obtain
\begin{align}
    &\sum_{\alpha \in \mathcal{F}_{D}} \lambda_{f_{\psi}}(N(\alpha)) \widehat{H_A}\left( \alpha,m,-\frac{\mathrm{Arg}(\alpha/\overline{\alpha})}{2 \pi} \right) \label{BoundInnerSumsSa-5} \\
    &\leq \sum_{\substack{\alpha \in \mathcal{F}_{D} \\ A/2 \leq N(\alpha) \leq 2A}} |\lambda_{f_{\psi}}(N(\alpha))| \cdot \left| \widehat{H_A}\left( \alpha,m,-\frac{\mathrm{Arg}(\alpha/\overline{\alpha})}{2 \pi} \right) \right| \nonumber \\
    &\ll_{D,\varepsilon, \delta} \ell^{\varepsilon} T^{1/2} (\ell A)^{-1/2} \sum_{\alpha\in S_{T^{-1+\delta}}} |\lambda_{f_{\psi}}(N(\alpha))| + \ell^{\varepsilon} T^{-3/2} (\ell A)^{-1/2} \sum_{\substack{\alpha \in \mathcal{F}_{D} \\ A/2 \leq N(\alpha) \leq 2A \\ \alpha \notin S_{T^{-1+\delta}} }} |\lambda_{f_{\psi}}(N(\alpha))| \nonumber \\
    &\ll_{D,\varepsilon, \delta} \ell^{\varepsilon} T^{1/2} (\ell A)^{-1/2} \sum_{\alpha\in S_{T^{-1+\delta}}} |\lambda_{f_{\psi}}(N(\alpha))| + \ell^{\varepsilon} T^{-3/2} (\ell A)^{-1/2} \sum_{\substack{\alpha \in \mathcal{F}_{D} \\ A/2 \leq N(\alpha) \leq 2A}} |\lambda_{f_{\psi}}(N(\alpha))|. \nonumber
\end{align}
The result follows from \eqref{BoundInnerSumsSa-1}, \eqref{BoundInnerSumsSa-2}, and \eqref{BoundInnerSumsSa-5}.
\end{proof}

\end{document}
